\section{Pérennité, maintenance et gestion des risques}
Le passage d'XTF vers un nouveau système ne se limite pas uniquement à un arbitrage d'ingénierie logicielle. Il constitue un défi de gouvernance et de gestion des ressources humaines et de viabilité du logiciel sur le long terme pour la bibliothèque Cujas.

Tout d'abord, à l'instar d'autres bibliothèques, l'ensemble de l'infrastructure technique repose principalement sur les compétences d'une seule personne. Ce constat d'une dépendance extrême est reconnu par le directeur et les experts métiers, attendant avec la nouvelle infrastructure d'avoir une meilleure visibilité dans le back-office. Pour éviter de créer à nouveau une \enquote{boîte noire} et d'abrégrer prématurément la durée de vie de la bibliothèque numérique, toute refonte doit alors s'accompagner d'une documentation fonctionnelle et technique très détaillée. Cette documentation permet ainsi de comprendre et de justifier des choix qui ont été fait. Elle doit consigner \enquote{chaque choix de circuit, l'architecture des bases de données, le code des API et la configuration des serveurs} permettant ainsi de disposer d'une compréhension globale dans le cas du départ du concepteur, remplacé par une autre personne assurant sa pérennité.
	
Le choix de pérenniser un prototype sur-mesure amène la bibliothèque Cujas à devoir assurer la totalité de \enquote{l'effort de maintenance corrective en interne, de mise à jour des sécurités et à l'adaptation à de futurs standards}. Or, la bibliothèque Cujas ne dispose pas d'un vivier de développeurs permettant d'assurer un suivi de la bibliothèque numérique. Contrairement à la maitenance d'un système open-source permet de partager l'effort réduisant la charge technique.

La refonte de la bibliothèque numérique de la bibliothèque cujas ne saurait se résumer uniquement à un simple chantier informatique. Les choix d'infrastructure posent les fondations matérielles et logicielles indispensables à la pérennité, à la sécurité et à l'interopérabilité des collections patrimoniales. Cependant, aussi évolué que peut être l'ingénierie logicielle, elle ne serait être validée si elle ne s'adapte pas aux pratiques contemporaines de la recherche. Il est donc nécessaire d'évaluer comment mettre en place des services fluides adaptés à l'usage.